% GDZ PPN236018647,native528–530;HKUST529–531.
\clearpage\SourceLeaf{528}\ArithmeticHead{499}{R}{1796 IUL.—1796 AUG. 18.}
\NumberedItem{23}
\indent Rationem theorematis aurei quomodo profundius perscrutari oporteat perspexi et ad hoc accingor supra quadraticas aequationes egredi conatus. Inventio formularum, quae semper per primos: \({}^{n}\!\surd1\) (numerice) dividi possunt.
\par\noindent\hfill[1796] Aug. 13.\quad Ibid. [Gottingae]

\medskip\begingroup\small
\indent Siehe Werke II, S. 230. Gemeint sind die Ausdrücke \(x^\pi-1\), für die \(p^n\equiv1\ (\mathrm{mod.}\,\pi)\). — Im letzten Satze des Textes hat die Handschrift \textls[120]{qui} statt \textls[120]{quae}.
\par\noindent\hfill\textsc{Bachmann.}
\par\endgroup

\NumberedItem{24}
\indent Obiter \((a+b\sqrt{-1})^{m+n\sqrt{-1}}\) evolutum.
\par\noindent\hfill[1796 Aug.] 14.

\medskip\begingroup\small
\indent Bei S. 111 des \textsc{Leiste} steht eine Aufzeichnung: »\textit{Canon quantitatum imaginariarum exponentialium}« mit Ausdrücken für \(a^{\surd-1}\) und \((a+b\surd-1)^{n\surd-1}\).
\par\noindent\hfill\textsc{Schlesinger.}
\par\endgroup

\NumberedItem{25}
\indent Rei summa iamiam intellecta. Restat ut singula muniantur.
\par\noindent\hfill[1796 Aug.] 16.\quad G[ottingae]

\medskip\begingroup\small
\indent Da diese Aufzeichnung, ebenso wie die Nummern 22, 23, 26, 27, in der Handschrift rot (in der Nachbildung im Halbton) unterstrichen ist, bezieht sie sich wahrscheinlich auf einen der in den genannten Nummern behandelten arithmetischen Gegenstände, während die Nr. 24 nur eine gelegentlich (obiter) gemachte Bemerkung darstellt.
\par\noindent\hfill\textsc{Schlesinger.}
\par\endgroup

\NumberedItem{26}
\indent \((a^p)\equiv(a)\ \mathrm{mod.}\,p,\ a\) radix aequationis cuiusvis quomodocunque irrationalis.
\par\noindent\hfill[1796 Aug.] 18.\quad [Gottingae]

\medskip\begingroup\small
\indent Es ist dies der Satz, der im art. 350. der \textit{Analysis Residuorum}, Werke II, S. 224, in der Form
\[
(P,\rho^\tau)\equiv P\ (\mathrm{mod.}\,\tau)
\]
\noindent ausgesprochen ist. Es bedeutet also hier \((a)\) eine ganze rationale Funktion von \(x\), die für \(x=a\) verschwindet, und entsprechend \((a^p)\) die Funktion, deren Wurzeln die \(p\)-ten Potenzen der Wurzeln der ersten Funktion sind; \(p\) Primzahl.
\par\noindent\hfill\textsc{Dedekind. Bachmann.}
\par\endgroup
\par\smallskip\noindent\hfill 63\textsuperscript{*}

\clearpage\SourceLeaf{529}\ArithmeticHead{500}{L}{TAGEBUCH MIT ERLÄUTERUNGEN.}
\NumberedItem{27}
\indent Si \(P,Q\) functiones alg[ebraicae] quantitatis indeterminatae fuerint inc[ognitae]. Datur:
\[
tP+uQ=1
\]
\noindent in algebra tum speciata tum numerica.
\par\noindent\hfill[1796 Aug.] 19.\quad G[ottingae]

\medskip\begingroup\small
\indent Siehe für die algebra numerica (Zahlentheorie) \textit{Analysis Residuorum}, art. 335., Werke II, S. 215, für die algebra speciata oder speciosa (Buchstabenrechnung) \textit{Demonstratio nova etc.}, 1816, art. 2., Werke III, S. 35.
\par\noindent\hfill\textsc{Klein. Bachmann.}
\par\endgroup

\NumberedItem{28}
\indent Exprimuntur potestates radicum aequationis propositae aggregatae per coefficientes aequationis lege perquam simplici (cum aliis quibusdam geometr[icis] in Exerc[itationibus]).
\par\noindent\hfill[1796 Aug.] 21.\quad G[ottingae]

\medskip\begingroup\small
\indent Vergl. die Bemerkung zu der Nr. 6, oben S. 490. Der art. [2.] der dort angeführten Leisteaufzeichnung (oben S. 128) zeigt, daß es sich hier um die \textls[120]{explizite} Darstellung der Potenzsummen der Wurzeln durch die Koeffizienten einer algebraischen Gleichung handelt. Für \(n=1,2,3,4\) finden sich diese Formeln bei \textsc{Albert Girard} in der \textit{Invention nouvelle en l’Algèbre}, Amsterdam 1629 (neue Ausgabe Leyden 1884), die allgemeine Darstellung mit dem von \textsc{Gauss} gefundenen Bildungsgesetz gibt \textsc{Waring} in den \textit{Miscellanea analytica de aequationibus algebraicis etc.}, Cambridge 1762\footnote{Vergl. hierzu L. \textsc{Saalschütz}, Bibliotheca mathematica (3) \textbf{9}, 1908, S. 65.}. In \textsc{Eulers} \textit{Introductio in analysin infinitorum}, Lausannae 1748, Liber I, § 166, S. 128 findet man bloss die \textsc{Newton}schen Formeln, die keine independente, sondern nur eine rekurrente Darstellung der Potenzsummen geben. In der Abhandlung \textit{Observationes circa radices aequationum}\begingroup\renewcommand{\thefootnote}{**)}\footnote{Novi Comm. Acad. Petrop. \textbf{15} (1770) 1771, S. 51.}\endgroup\ hat auch \textsc{Euler} eine independente Darstellung der Potenzsummen abzuleiten versucht, doch ist er zu keinem so übersichtlichen Bildungsgesetz gelangt, wie es bei \textsc{Waring} und \textsc{Gauss} vorliegt. Wie wenig bekannt \textsc{Warings} Darstellung geblieben war, geht u. a. daraus hervor, daß von ihr in \textsc{Klügels} \textit{Mathematischem Wörterbuch} I, Leipzig 1803, wo S. 467 und 507 (fälschlich als 495 numeriert) von den \textsc{Newton}schen und den expliziten \textsc{Girard}schen Formeln gehandelt wird, nicht die Rede ist. Auch die im art. [1.] der Leisteaufzeichnung (oben S. 127) angegebene explizite Darstellung der Gleichungskoeffizienten durch die Potenzsummen findet sich bei \textsc{Waring} a. a. O., und zwar gibt \textsc{Waring} die allgemeinen Formeln für ein beliebiges \(n\), während \textsc{Gauss} sich auf die vier ersten Koeffizienten beschränkt.

\indent Die in unserer Tagebuchnotiz erwähnten \textit{Exercitationes Mathematicae}, die oben S. 138 abgedruckt sind, haben bei den artt. [1.], [2.] in der Tat dieselbe Zeitangabe, 21. August 1796.
\par\noindent\hfill\textsc{Loewy.}
\par\endgroup

\clearpage\SourceLeaf{530}\ArithmeticHead{501}{R}{1796 AUG. 19—1796 SEPT. 2.}
\NumberedItem{29}
\indent Summatio Seriei infinitae
\[
1+\frac{x^n}{1\ldotp\,\ldotp\,\ldotp n}+\frac{x^{2n}}{1\ldotp\,\ldotp\,\ldotp2n}\ \text{etc.}
\]
\noindent\hfill eod[em die, 1796 Aug. 21.]

\medskip\begingroup\small
\indent Die Reihe befriedigt die Differentialgleichung
\[
\frac{d^n u}{dx^n}=u
\]
\noindent mit den Anfangsbedingungen, daß für \(x=0,\ u=1\) und die \(n-1\) ersten Derivierten gleich Null sind. Daraus ergibt sich
\[
\sum_{\nu=0}^{\infty}\frac{x^{\nu n}}{(\nu n)!}=\frac1n\sum_{\kappa=1}^{n}e^{x\left(\cos\frac{2\kappa\pi}{n}+i\sin\frac{2\kappa\pi}{n}\right)};
\]
\noindent vergl. die Bemerkung zum art. [3.] der \textit{Exercitationes mathematicae}, oben S. 143, 144. Über Funktionen dieser Art, die \textsc{Hoene-Wronski} (\textit{Philosophie des Mathématiques}, Paris 1811) als \textls[120]{Sinus} bezw. \textls[120]{Kosinus höherer Ordnung} bezeichnet hat, gibt es eine umfangreiche Literatur, die bei S. \textsc{Günther}, \textit{Die Lehre von den . . . Hyperbelfunktionen}, Halle 1881, und bei K. \textsc{Wallner}, Programm der K. Realschule Rothenburg o. T. 1913, zusammengestellt ist.
\par\noindent\hfill\textsc{Schlesinger.}
\par\endgroup

\par\medskip\Needspace{4\baselineskip}\begin{center}[30].\end{center}
\indent Minutiis quibusdam exceptis feliciter scopum attigi scil[icet] si
\[
p^n\equiv1\ (\mathrm{mod}\,\pi),
\]
\noindent fore \(x^\pi-1\) compositum e factoribus gradum \(n\) non excedentibus \textls[120]{et proin aequationem conditionalem fore solubilem}; unde duas theor[ematis] aurei demonstr[ationes] deduxi.
\par\noindent\hfill[1796] Sept. 2.\quad G[ottingae]

\medskip\begingroup\small
\indent Siehe \textit{Analysis Residuorum} art. 360., Werke II, S. 230, sowie die Tagebuchnotizen Nr. 23 vom 13. August 1796 und Nr. 68 vom 21. Juli 1797. Die \textls[120]{aequatio conditionalis} ist vermutlich dieselbe, die \textsc{Gauss} sonst \textls[120]{aequatio auxiliaris} genannt hat, vergl. \textit{Analysis Residuorum} artt. 365. und 366., Werke II, S. 233, 234, wo sich die beiden Beweise des »theorematis aurei« finden. Die \textls[120]{Minutia excepta} sind die am Ende des art. 363., Werke II, S. 232 erwähnten Schwierigkeiten, die erst in der Tagebuchnotiz Nr. 68 als behoben angemerkt werden.
\par\noindent\hfill\textsc{Dedekind. Bachmann.}
\par\endgroup
